\begin{problem}{Zero Segment}{standard input}{standard output}{2 seconds}{256 megabytes}

You are given an array $a$ of $n$ positive integers, and an integer $k$.

You may perform the following operations on the array any number of times, in any order:

\begin{itemize}
  \item \textbf{Pay:} choose an index $i$ ($1 \le i \le n$), pay $a_i$ coins, and set $a_i := 0$.
  \item \textbf{Free:} choose three indices $l$, $m$, $r$ such that $1 \le l < m < r \le n$, and set $a_m := \max(a_l, a_r)$. This operation costs nothing.
\end{itemize}

Find the minimum total number of coins you have to pay so that some contiguous segment of length at least $k$ of the array consists only of zeros.


\InputFile
The first line contains a single integer $t$ ($1 \le t \le 2 \cdot 10^5$) --- the number of test cases.

The first line of each test case contains two integers $n$ and $k$ ($1 \le k \le n \le 2 \cdot 10^5$) --- the size of the array and the minimum length of the required zero segment.

The second line of each test case contains $n$ integers $a_1, a_2, \ldots, a_n$ ($1 \le a_i \le 10^9$) --- the elements of the array.

It is guaranteed that the sum of $n$ over all test cases does not exceed $2 \cdot 10^5$.


\OutputFile
For each test case, print a single integer --- the minimum total cost needed to obtain a contiguous zero segment of length at least $k$.


\Examples

\begin{example}
\exmpfile{example.01}{example.01.a}%
\exmpfile{example.02}{example.02.a}%
\end{example}

\Note
In the first test case of the first example, one way to spend exactly $3$ coins is as follows. Pay $1$ coin to set $a_1 := 0$. Apply the free operation with $(l, m, r) = (1, 2, 5)$: this sets $a_2 := \max(a_1, a_5) = \max(0, 5) = 5$ --- note that this does not help at all. Now pay $2$ coins to set $a_3 := 0$, and apply the free operation with $(l, m, r) = (1, 2, 3)$ to set $a_2 := \max(a_1, a_3) = 0$. The segment $[1, 3]$ is all zeros. Trying to zero a longer segment, say $[2, 4]$, turns out to cost more, and careful consideration of every possibility shows that cost $3$ cannot be beaten.

In the second test case of the first example the answer is $40$: a suitable sequence of operations achieves it, and one can check that no cheaper sequence exists.

In the third test case of the first example, $k = 1$, so after paying $7$ coins for the only element the array already contains a zero segment of length $1$.

In the second example, for the first test case the cheapest single payment has cost $1$, which is sufficient. For the second test case, the optimal cost is $9$.


\end{problem}

